\chapter{Происхождение общего теплового следа из Real-суперсвязности}
\label{ch:v7-real-superconnection-common-trace-origin-gate}

\section{Решающий вопрос}

Предыдущий гейт получил положительный локальный гессиан формального
функционала
\begin{equation}
 \mathcal S_{\rm form}
 =-\Tr_{H_{15}}e^{-\mathfrak m_{\rm ch}^2}
  -\Tr_Ee^{-\mathfrak m_C^2}
  +\Tr_{\rm phys}e^{-\Phi^2}.
 \label{eq:v7-real-trace-formal-functional}
\end{equation}
Проверим, могут ли два минуса и один плюс быть обычной градуировкой одной
конечномерной Real-суперсвязности, а не назначением ролей после вычисления.

\section{Индексный запрет}

Пусть $\mathcal H=\mathcal H^+\oplus\mathcal H^-$ и нуль-форменная часть
суперсвязности нечётна:
\begin{equation}
 \mathbb A_0=
 \begin{pmatrix}0&B^*\\B&0\end{pmatrix},
 \qquad
 \mathbb A_0^2=
 \begin{pmatrix}B^*B&0\\0&BB^*\end{pmatrix}.
 \label{eq:v7-real-trace-odd-square}
\end{equation}
Ненулевые спектры $B^*B$ и $BB^*$ совпадают с кратностями. Поэтому для
любого $t>0$
\begin{equation}
 \Str e^{-t\mathbb A_0^2}
 =\dim\ker B-\dim\ker B^*
 =\dim\mathcal H^+-\dim\mathcal H^-.
 \label{eq:v7-real-trace-index-identity}
\end{equation}
Это конечномерная форма механизма Маккина--Зингера. Она постоянна по $B$;
все её корневые производные и гессиан равны нулю. Прямая сумма трёх таких
комплексов даёт только сумму индексов:
\begin{equation}
 \Str e^{-t(\mathbb A_{0,1}\oplus\mathbb A_{0,2}\oplus
 \mathbb A_{0,3})^2}=\operatorname{ind}B_1+
 \operatorname{ind}B_2+\operatorname{ind}B_3.
 \label{eq:v7-real-trace-direct-sum-index}
\end{equation}
Следовательно, она не может совпасть с \eqref{eq:v7-real-trace-formal-functional},
чей полный гессиан имеет минимальную моду $1.14399252085$.

\section{Почему ролевой знак не является решением}

На формальной прямой сумме можно ввести
\begin{equation}
 K_{\rm role}=\operatorname{diag}(-I_{H_{15}},-I_E,I_{\rm phys}),
 \qquad
 \Tr K_{\rm role}e^{-Q^2}=\mathcal S_{\rm form}.
 \label{eq:v7-real-trace-role-involution}
\end{equation}
Но $Q=\operatorname{diag}(\mathfrak m_{\rm ch},\mathfrak m_C,\Phi)$
коммутирует с $K_{\rm role}$ и не является нечётным относительно него:
\begin{equation}
 \{K_{\rm role},Q\}=2K_{\rm role}Q\ne0
 \quad\text{при}\quad Q\ne0.
 \label{eq:v7-real-trace-role-oddness-defect}
\end{equation}
Если исправить нечётность каноническим удвоением
\begin{equation}
 D_Q=\begin{pmatrix}0&Q^*\\Q&0\end{pmatrix},
 \qquad \Str e^{-tD_Q^2}=0
 \label{eq:v7-real-trace-canonical-doubling}
\end{equation}
для квадратного $Q$. Real-удвоение и последующий физический полуслед лишь
удваивают и затем делят это сокращение; нужная динамика не возвращается.

\begin{theorem}[No-go обычного теплового суперследа]
Функционал \eqref{eq:v7-real-trace-formal-functional} не является
нуль-форменной компонентой обычного теплового суперследа одной
конечномерной нечётной Real-суперсвязности. Честный суперслед сокращается до
индекса, а ролевой оператор $K_{\rm role}$ воспроизводит знаки только как
дополнительная вставка, не как градуировка данного ненулевого $Q$.
\end{theorem}

\section{Точная граница результата}

Запрет не относится к высшим дифференциально-форменным компонентам характера
Черна, трансгрессиям, условным или эквивариантным следам и положительной
норме кривизны. Но каждый такой маршрут содержит дополнительную
геометрическую структуру и не может называться обычным суперследовым выводом
\eqref{eq:v7-real-trace-formal-functional}. Литература Квиллена и
Бисмута--Фрида подтверждает характерологическую и трансгрессионную роль
суперследа, но не предоставляет требуемый скалярный потенциал автоматически
\cite{v7-real-trace-quillen,v7-real-trace-bismut-freed}.

\begin{nogocriterion}[Запрет ручной ролевой градуировки]
Нельзя объявлять $K_{\rm role}$ физически выведенным только потому, что его
вставка сохраняет найденный положительный гессиан. Требуется получить
относительную инволюцию из уже существующего дифференциала, Real-структуры,
условного ожидания или одной положительной нормы кривизны.
\end{nogocriterion}

\begin{protocol}
\begin{enumerate}
 \item обычный нечётный тепловой суперслед: \textbf{индексный};
 \item зависимость его нуль-форменной части от полей: \textbf{нет};
 \item прямая сумма трёх комплексов: \textbf{не помогает};
 \item Real-удвоение и полуслед: \textbf{не помогают};
 \item ролевой оператор воспроизводит знаки: \textbf{формально да};
 \item происхождение ролевого оператора: \textbf{не выведено};
 \item формальный локальный вакуум предыдущего гейта: \textbf{сохраняется};
 \item его повышение до физического родителя: \textbf{не пройдено};
 \item следующий гейт: \textbf{производная относительная инволюция или
 норма общей кривизны}.
\end{enumerate}
\end{protocol}

Машинный сертификат сохранён в
\path{s2t_v7_real_superconnection_common_trace_origin_gate_results.json}.

\begin{thebibliography}{9}
\bibitem{v7-real-trace-quillen}
D.~Quillen,
\emph{Superconnections and the Chern Character},
Topology 24 (1985), 89--95.
\bibitem{v7-real-trace-mckean-singer}
H.~P.~McKean, I.~M.~Singer,
\emph{Curvature and the Eigenvalues of the Laplacian},
J. Differential Geometry 1 (1967), 43--69.
\bibitem{v7-real-trace-bismut-freed}
J.-M.~Bismut, D.~S.~Freed,
\emph{The Analysis of Elliptic Families. I. Metrics and Connections on
Determinant Bundles},
Comm. Math. Phys. 106 (1986), 159--176.
\end{thebibliography}